As briefly explained in \cref{bs:range}, a target with a radial speed with respect to the radar will produce a frequency shift on the reflected signal called the Doppler shift. This frequency shift is due to the so-called bulk motion of the object, that is the general velocity at which an object is moving. For many types of targets, however, oscillatory motions, which can also be called \textit{micro motions}, might be present. Examples of micro motions include a rotating propeller of an aircraft, a rotation antenna or the flapping wings of bird. For our scenario micro motions are found in the swinging of arms an legs of a pedestrians or the movement of the legs of a cyclist. For vehicles, the rotation of wheels could be considered a micro motion.

These micro motions cause for a further frequency shift in the transmitted signals called the \textit{Micro-Doppler effect}\cite{chen_micro-doppler_2011}. For periodic vibrations or rotations, such as the case of the limbs of a pedestrian, shifts about the the center of the Doppler-shift due to the bulk motion are observed. 

\begin{figure}[h!]
	\centering
	\begin{tikzpicture}
% Reference Grid
\draw[gray,very thin] (0,0) grid (10,4); 

\end{tikzpicture}
	\caption{Result to illustrate the Micro Doppler Effect}
\end{figure} 